\documentclass[11pt]{article}

% LuaLaTeX
\usepackage{fontspec}
\usepackage{unicode-math}
\setmainfont{Latin Modern Roman}
\setmathfont{Latin Modern Math}

\usepackage[margin=1in]{geometry}
\usepackage{setspace}
\setstretch{1.15}

\usepackage{amsmath,amsthm,mathtools}
\usepackage{microtype}
\usepackage{booktabs}
\usepackage{hyperref}
\hypersetup{
  colorlinks=true,
  linkcolor=blue,
  citecolor=blue,
  urlcolor=blue
}

\newtheorem{theorem}{Theorem}[section]
\newtheorem{lemma}[theorem]{Lemma}
\newtheorem{corollary}[theorem]{Corollary}
\newtheorem{proposition}[theorem]{Proposition}
\newtheorem{claim}[theorem]{Claim}
\theoremstyle{definition}
\newtheorem{definition}[theorem]{Definition}
\theoremstyle{remark}
\newtheorem{remark}[theorem]{Remark}

\newcommand{\histories}{\mathcal{H}}
\newcommand{\valuation}{V}
\newcommand{\worlds}{W}
\newcommand{\acts}{\mathcal{A}}
\newcommand{\obs}{\mathcal{O}}
\newcommand{\admissible}{\mathrm{Adm}}
\newcommand{\manifold}{\mathcal{M}}

\title{\textbf{The Central Theorem: Epistemic Overlap}\\
Admissible Continuation, Information, and Local-to-Global Consistency}
\author{Flyxion}
\date{October 2026}

\begin{document}
\maketitle

\begin{abstract}
This paper develops a common formal structure for several problems involving
admissible continuation under incomplete information. The central result,
Epistemic Overlap (CT-000), characterizes when a deterministic policy based
only on observations can guarantee an admissible action. Such a policy exists
precisely when every observational equivalence class has a nonempty common
admissible-action set.

The result is elementary but useful because it separates two questions:
whether each possible world admits some action, and whether the information
available to an agent is sufficient to select an action admissible in every
world compatible with that information.

Subsequent sections relate this structure to admissible histories,
local-to-global geometric models, bounded-correction models, minimal
decision-relevant history, and durable adjudication. These later
correspondences have different logical status. Standard mathematical results,
conditional propositions, modeling assumptions, and implementation principles
are distinguished rather than treated as interchangeable forms of proof.

Stable identifiers CT-000 through CT-016 correspond to the claim registry in
\texttt{SPEC.md}.
\end{abstract}

% =====================================================================
\section{Specification}


The notation, claim identifiers, status terminology, and scope conventions used throughout the paper are defined below. Mathematical results are distinguished from statements conditional on model assumptions and from requirements pertaining to implementation.

\subsection{Claim identifiers}

Claims are assigned stable identifiers of the form
\[
  \texttt{CT-000},\texttt{ CT-001},\ldots,\texttt{ CT-016}.
\]
The identifier denotes a claim, not a particular proof, implementation,
formalization, file, or revision. Multiple artifacts may therefore refer to
the same identifier without being treated as equivalent evidence for it.

The canonical registry is maintained in \texttt{SPEC.md}. The present paper
contains the mathematical statements and their qualifications.

\subsection{Evidence and status}

The following distinctions are used throughout.

\begin{description}
  \item[Proved.]
  A mathematical statement is established from explicitly stated assumptions.

  \item[Conditional.]
  A mathematical conclusion follows within a specified model or under
  additional hypotheses that are not asserted universally.

  \item[Standard result.]
  The mathematical statement is a standard result used as part of the model;
  its interpretation in the present setting is separate from the theorem
  itself.

  \item[Modeling correspondence.]
  A mathematical structure is proposed as a useful representation of some
  computational or semantic structure. The correspondence is not itself
  asserted as a mathematical theorem.

  \item[Qualitative.]
  The claim is motivated by the source model but has not been reduced to a
  sufficiently explicit mathematical statement for a general proof.

  \item[Implementation principle.]
  The claim constrains or motivates system design. An implementation can
  satisfy or test the constraint, but doing so is not a proof of a general
  mathematical proposition.

  \item[Formalized.]
  A stated proposition has been represented and checked in a proof assistant.
  Formalization establishes the encoded proposition relative to its
  definitions and assumptions; it does not automatically establish broader
  interpretations attached to the claim.

  \item[Implemented.]
  Software exhibits behavior corresponding to a claim or design principle.
  Implementation correspondence is distinct from mathematical proof.
\end{description}

These categories may overlap. For example, a proved mathematical statement
may also have a machine-checked formalization and an implementation that
satisfies a corresponding invariant.

\subsection{Claim registry}

\begin{center}
\small
\begin{tabular}{@{}p{1.4cm}p{7.3cm}p{4.7cm}@{}}
\toprule
ID & Claim & Status \\
\midrule
CT-000 &
Epistemic Overlap and the No-Safe-Action corollary &
proved \\

CT-001 &
Admissibility is independent of valuation when admissibility is defined
structurally &
established within the stated model \\

CT-002 &
Prefix closure of admissible histories &
established under the prefix-closure assumption \\

CT-003 &
Unique continuation under a committed prefix &
established within the stated model \\

CT-004 &
Persistence of disfavored admissible trajectories &
qualitative \\

CT-005 &
Pruning excludes continuations rather than repairing prior states &
established within the stated model \\

CT-006 &
Tangent-normal decomposition and its use as a model of admissible
transformation &
standard result plus modeling correspondence \\

CT-007 &
Off-manifold generation may enlarge effective output support under additional
regularity assumptions &
conditional \\

CT-008 &
Local-chart coherence and tangent-flow interpretation &
modeling correspondence \\

CT-009 &
Local admissibility does not by itself imply global coherence &
general structural observation; geometric form model-dependent \\

CT-010 &
Bounded correction cannot uniformly offset the modeled production term under
A1--A4 &
conditional \\

CT-011 &
Finite safe regime under bounded protective capacity and the additional
monotonicity assumptions &
conditional \\

CT-012 &
Increasing marginal risk under the specified visibility scaling model &
model-dependent \\

CT-013 &
Local stress criterion for semantic stability &
model-dependent \\

CT-014 &
Scaling formulation of visibility-induced instability &
model-dependent \\

CT-015 &
Admissibility may factor through a reduced decision-relevant history &
proved as a factorization result \\

CT-016 &
Adjudications affecting future admissibility must survive authoritative
reconstruction &
implementation principle \\
\bottomrule
\end{tabular}
\end{center}

\subsection{Source correspondence}

The claims consolidate material developed in several related texts. The
correspondence is:

\begin{center}
\small
\begin{tabular}{@{}p{5.3cm}p{5.0cm}p{2.7cm}@{}}
\toprule
Source & Relevant material & Claims \\
\midrule
\texttt{Admissible Histories.tex} &
History, prefix, continuation, pruning &
CT-001--CT-005 \\

\texttt{Never-Predict-Noise.tex} &
Tangent-normal decomposition, off-manifold analysis, local coherence,
curvature &
CT-006--CT-009 \\

\texttt{Against\_Indiscriminate\_Visibility.tex} &
Visibility, bounded correction, stress, scaling &
CT-010--CT-014 \\

\texttt{resources/admissible-continuation.tex} &
Admissible continuation &
CT-000, CT-003 \\

\texttt{physics/cosmology/admissible-geometry.tex} &
Geometric formulation of admissibility &
CT-006 \\
\bottomrule
\end{tabular}
\end{center}

These correspondences identify intellectual provenance within the corpus.
They do not imply that every source uses exactly the definitions adopted in
this consolidated treatment.

\subsection{Formal structure of CT-000}

The central theorem uses the following objects:

\begin{itemize}
  \item a set of possible worlds $\worlds$;
  \item a set of observations $\obs$;
  \item a set of actions $\acts$;
  \item an observation map
    \[
      O:\worlds\to\obs;
    \]
  \item an admissible-action map
    \[
      \admissible:\worlds\to\mathcal{P}(\acts);
    \]
  \item a deterministic observation-based policy
    \[
      \pi:\obs\to\acts.
    \]
\end{itemize}

For $o\in O(\worlds)$, the observation class is
\[
  [o]=\{w\in\worlds:O(w)=o\}.
\]

The common admissible-action set associated with $o$ is
\[
  C(o)
  =
  \bigcap_{w\in[o]}\admissible(w).
\]

The theorem concerns the equivalence between

\[
  \exists\pi:\obs\to\acts
  \quad
  \forall w\in\worlds,\;
  \pi(O(w))\in\admissible(w)
\]
and
\[
  \forall o\in O(\worlds),\;
  C(o)\neq\varnothing,
\]
subject to the totality and selection conditions stated below.

This is stronger than requiring pairwise overlap among admissible-action
sets. In general,
\[
  \admissible(w_i)\cap\admissible(w_j)\neq\varnothing
  \qquad\text{for every }i,j
\]
does not imply
\[
  \bigcap_i\admissible(w_i)\neq\varnothing.
\]
The theorem therefore uses the intersection of the entire observation class.

\subsection{Totality and choice}

Two logically distinct issues arise when constructing a policy.

First, the paper defines a policy on all of $\obs$, although only values in
$O(\worlds)$ are used to determine actions in possible worlds. If
\[
  \obs\setminus O(\worlds)\neq\varnothing,
\]
then a total policy $\pi:\obs\to\acts$ must assign actions to unrealized
observations as well. Thus $\acts$ must be nonempty, or the policy domain must
instead be restricted to $O(\worlds)$.

Second, the forward construction selects one action from every nonempty set
$C(o)$ defined above.
A simultaneous selection from an arbitrary family of such sets requires an
appropriate choice principle or an explicit selection rule.

For finitely many relevant observation classes, finite choice suffices.
If the action sets carry an ordering or another rule that canonically selects
an element, that rule can replace abstract choice.

Restricting the policy to
\[
  \pi:O(\worlds)\to\acts
\]
eliminates the need to assign arbitrary actions to unrealized observations.
It does not, in general, eliminate the selection problem among the nonempty
sets $C(o)$.

Accordingly, two versions of CT-000 may be distinguished:

\begin{description}
  \item[Total-policy version.]
  $\pi:\obs\to\acts$, requiring values outside $O(\worlds)$ when such
  observations exist.

  \item[Realized-observation version.]
  $\pi:O(\worlds)\to\acts$, requiring values only for observations actually
  produced by worlds in $\worlds$.
\end{description}

The substantive overlap condition is identical in both versions.

\subsection{Deterministic scope}

CT-000 concerns deterministic policies. Randomization does not automatically
remove the obstruction identified by the theorem.

If admissibility requires the selected action itself to belong to
$\admissible(w)$ with certainty, then a randomized policy supported on actions
outside the common intersection fails in at least one compatible world,
provided the observation class is countable. The uncountable case is treated
in the subsection on randomized policies.
A randomized analogue would therefore require an explicitly defined
probabilistic admissibility criterion.

No such criterion is assumed here.

Likewise, CT-000 is a static existence result. It does not by itself specify
learning, exploration, belief revision, communication between agents, or
sequential information acquisition. Those extensions require additional
structure.

\subsection{Admissibility}

The abstract theorem places no intrinsic restriction on the origin of
$\admissible(w)$. It may represent legal actions, continuation constraints,
safety conditions, protocol states, physical feasibility, or another
world-dependent admissibility relation.

In the applications developed in this paper, admissibility is usually
interpreted structurally: an action is admissible when it constitutes a legal
continuation of the relevant attested history.

That interpretation belongs to the admissible-history model. CT-000 itself
requires only the set-valued map
\[
  \admissible:\worlds\to\mathcal{P}(\acts).
\]

This distinction matters because the central theorem does not depend on
CT-001--CT-016 for its validity. Those claims provide applications,
interpretations, and related structures rather than premises required to
prove CT-000.

\subsection{Information partitions}

The observation map $O$ induces an equivalence relation
\[
  w\sim_O w'
  \quad\Longleftrightarrow\quad
  O(w)=O(w').
\]

The equivalence classes of $\sim_O$ are precisely the information sets on
which an observation-based deterministic policy must be constant.

CT-000 can therefore be stated without explicit reference to observations:
a policy constant on each information class is globally admissible if and
only if each information class has a common admissible action.

This formulation makes clear that the theorem concerns information loss
induced by a partition of the world space. The particular representation of
the observations is secondary.

\subsection{Refinement of information}

Suppose two observation maps
\[
  O_1:\worlds\to\obs_1,
  \qquad
  O_2:\worlds\to\obs_2
\]
satisfy
\[
  O_1=f\circ O_2
\]
for some function $f$. Then $O_2$ is at least as informative as $O_1$:
every $O_2$-class is contained in an $O_1$-class.

This immediately yields a useful monotonicity property.

\begin{proposition}[Information refinement]
If a globally admissible policy exists under $O_1$, then a globally
admissible policy exists under any refinement $O_2$ of $O_1$.
\end{proposition}

\begin{proof}
Every $O_2$-class is a subset of an $O_1$-class. If
\[
  \bigcap_{w:O_1(w)=o_1}\admissible(w)
  \neq\varnothing,
\]
then removing worlds from that intersection cannot make it smaller.
Consequently every refined class also has a nonempty common admissible-action
set. CT-000 then supplies a globally admissible policy for $O_2$.
\end{proof}

\begin{remark}
The converse does not hold. Coarsening observations can merge classes whose
admissible-action sets have no common element.

Thus, in the abstract CT-000 model, additional information cannot destroy the
\emph{existence} of an admissible policy. Claims elsewhere in the paper about
costs associated with increased visibility concern additional mechanisms and
must not be identified with refinement of $O$ alone.
\end{remark}

\subsection{Refinement order}

Observation maps can be partially ordered by informational refinement. Write
\[
  O_2 \succeq O_1
\]
when there exists a function $f$ such that
\[
  O_1=f\circ O_2.
\]
Equivalently, every equivalence class induced by $O_2$ is contained in an
equivalence class induced by $O_1$.

Call an observation map \emph{admissibility-sufficient} when it admits a
globally admissible observation-based policy.

\begin{proposition}[Upward closure under refinement]
The class of admissibility-sufficient observation maps is upward closed under
informational refinement. If
\[
  O_2\succeq O_1
\]
and $O_1$ is admissibility-sufficient, then $O_2$ is
admissibility-sufficient.
\end{proposition}

\begin{proof}
Let
\[
  O_1=f\circ O_2
\]
and let
\[
  \pi_1:\obs_1\to\acts
\]
be globally admissible under $O_1$. Define
\[
  \pi_2=\pi_1\circ f.
\]
Then for every $w\in\worlds$,
\[
  \pi_2(O_2(w))
  =
  \pi_1(f(O_2(w)))
  =
  \pi_1(O_1(w))
  \in
  \admissible(w).
\]
Hence $\pi_2$ is globally admissible under $O_2$.
\end{proof}

\begin{corollary}[Common refinements]
If $O_1$ is admissibility-sufficient and $O_2$ is any observation map, then
the joint observation
\[
  O_{12}(w)
  =
  \bigl(O_1(w),O_2(w)\bigr)
\]
is admissibility-sufficient.
\end{corollary}

\begin{proof}
The projection
\[
  \operatorname{pr}_1:
  \obs_1\times\obs_2\to\obs_1
\]
satisfies
\[
  O_1
  =
  \operatorname{pr}_1\circ O_{12}.
\]
Thus $O_{12}$ refines $O_1$, and the result follows from refinement
monotonicity.
\end{proof}

\begin{remark}
Admissibility-sufficiency is not generally preserved by coarsening. A
coarsening may merge two previously distinguishable classes whose
admissible-action sets have empty common intersection.

Consequently, the admissibility-sufficient information structures form an
upper set in the refinement order, not in general a lower set.
\end{remark}

\subsection{Finite obstruction certificates}

When the action space is finite, failure of epistemic overlap always has a
finite witness whose size is bounded by the number of available actions.

\begin{theorem}[Finite obstruction certificate]
Suppose
\[
  |\acts|=m<\infty.
\]
If for some realized observation $o$
\[
  \bigcap_{w\in[o]}\admissible(w)
  =
  \varnothing,
\]
then there exist worlds
\[
  w_1,\ldots,w_k\in[o],
  \qquad
  k\leq m,
\]
such that
\[
  \bigcap_{i=1}^{k}\admissible(w_i)
  =
  \varnothing.
\]
\end{theorem}

\begin{proof}
For every action $a\in\acts$, the assumption
\[
  \bigcap_{w\in[o]}\admissible(w)=\varnothing
\]
implies that $a$ fails to belong to at least one admissible-action set in the
observation class. Choose a world $w_a\in[o]$ such that
\[
  a\notin\admissible(w_a).
\]
The family
\[
  \{w_a:a\in\acts\}
\]
contains at most $m$ distinct worlds.

Suppose, toward a contradiction, that some action $b$ belongs to every
$\admissible(w_a)$. In particular,
\[
  b\in\admissible(w_b),
\]
contradicting the choice of $w_b$. Therefore
\[
  \bigcap_{a\in\acts}\admissible(w_a)
  =
  \varnothing.
\]
Removing duplicate worlds gives the required collection of at most $m$
worlds.
\end{proof}

\begin{corollary}[Bounded witness to policy failure]
For a finite action space of cardinality $m$, nonexistence of a globally
admissible policy can always be certified within a single observation class
by at most $m$ worlds in that class.
\end{corollary}

\begin{remark}
This strengthens the two-world obstruction without replacing it. When
$m>2$, policy failure need not be witnessed by a pair of worlds. For example,
three admissible-action sets may intersect pairwise while having empty total
intersection. The theorem nevertheless guarantees a finite obstruction whose
size is controlled by the action space rather than by the number of possible
worlds.
\end{remark}

\subsection{Convex admissibility and Helly reduction}

A different finite-witness result follows when the action space is geometric.

Suppose
\[
  \acts=\mathbb{R}^d
\]
and every $\admissible(w)$ is convex.

\begin{theorem}[Helly reduction for finite observation classes]
Let $[o]$ be a finite observation class. Then
\[
  \bigcap_{w\in[o]}\admissible(w)\neq\varnothing
\]
if and only if every subfamily of at most $d+1$ admissible-action sets has
nonempty intersection.
\end{theorem}

\begin{proof}
The forward implication is immediate. For the converse, the family
\[
  \{\admissible(w):w\in[o]\}
\]
is a finite family of convex subsets of $\mathbb{R}^d$. If every subfamily of
at most $d+1$ members has nonempty intersection, Helly's theorem implies
\[
  \bigcap_{w\in[o]}\admissible(w)
  \neq
  \varnothing.
\]
\end{proof}

\begin{corollary}[Low-dimensional obstruction]
For convex admissible-action sets in $\mathbb{R}^d$, failure of CT-000 on a
finite observation class has a witness involving at most $d+1$ worlds.
\end{corollary}

\begin{remark}
The finite-action and convex-action results bound obstruction size for
different reasons. The former depends on the cardinality of the action space;
the latter depends on its geometric dimension.

Neither result is required by CT-000. They identify settings in which the
potentially large intersection appearing in the theorem can be tested through
bounded subfamilies.
\end{remark}

\begin{remark}[Infinite families]
The finiteness of the observation class is essential to the stated form of
Helly's theorem. For infinite families of convex sets, nonempty intersection
of every subfamily of at most $d+1$ members does not by itself yield a
nonempty total intersection: the sets $[n,\infty)\subset\mathbb{R}$,
$n\in\mathbb{N}$, have every finite subfamily intersecting and empty total
intersection.

If all the sets are closed, at least one is compact, and every subfamily of at
most $d+1$ members has nonempty intersection, then Helly's theorem gives the
finite-intersection property, and compactness yields a nonempty total
intersection. Any extension of the Helly reduction to infinite observation
classes must state such conditions explicitly.
\end{remark}

\subsection{Randomized policies}

Let $(\acts,\Sigma)$ be a measurable space, assume every $\admissible(w)$
belongs to $\Sigma$, and let $\Delta(\acts)$ denote the probability measures
on $(\acts,\Sigma)$. A randomized observation-based policy is a map
\[
  \rho:\obs\to\Delta(\acts).
\]

Two notions of almost-sure admissibility must be distinguished. A policy
$\rho$ is \emph{pointwise almost-surely admissible} when
\[
  \rho(O(w))\bigl(\admissible(w)\bigr)=1
  \qquad
  \text{for every }w\in\worlds.
\]
It is \emph{robustly almost-surely admissible} when, for every realized
observation $o$, the set $C(o)$ belongs to $\Sigma$ and
\[
  \rho(o)\bigl(C(o)\bigr)=1.
\]
Robust admissibility implies pointwise admissibility. The converse can fail,
as shown below.

\begin{proposition}[Robust admissibility]
A robustly almost-surely admissible randomized policy exists only if
$C(o)\neq\varnothing$ for every realized observation $o$. Conversely, if every
$C(o)$ is nonempty and measurable and a selection from the sets $C(o)$ is
available, then a robustly almost-surely admissible randomized policy exists.
\end{proposition}

\begin{proof}
If $\rho(o)(C(o))=1$, then $C(o)\neq\varnothing$, since every probability
measure assigns measure $0$ to the empty set. Conversely, choose $a_o\in C(o)$
and let $\rho(o)$ be the point mass at $a_o$; then $\rho(o)(C(o))=1$.
\end{proof}

\begin{proposition}[Countable observation classes]
Suppose every observation class $[o]$ is countable. Then pointwise
almost-sure admissibility implies robust almost-sure admissibility. In
particular, if $C(o)=\varnothing$ for some realized $o$, no pointwise
almost-surely admissible randomized policy exists.
\end{proposition}

\begin{proof}
Fix a realized observation $o$ and enumerate $[o]=\{w_1,w_2,\ldots\}$, the
finite case being included. Then $C(o)=\bigcap_i\admissible(w_i)$ is
measurable. Pointwise admissibility gives
$\rho(o)(\acts\setminus\admissible(w_i))=0$ for every $i$. By countable
subadditivity,
\[
  \rho(o)\Bigl(\bigcup_i\bigl(\acts\setminus\admissible(w_i)\bigr)\Bigr)=0,
\]
and taking complements gives $\rho(o)(C(o))=1$.
\end{proof}

\begin{remark}[Uncountable classes]
The countability hypothesis cannot be dropped. Let $\acts=\worlds=[0,1]$
with a single observation and $\admissible(w)=[0,1]\setminus\{w\}$. Lebesgue
measure $\mu$ satisfies $\mu(\admissible(w))=1$ for every $w$, so the constant
policy $\rho(o)=\mu$ is pointwise almost-surely admissible. Nevertheless
$\bigcap_{w\in[0,1]}\admissible(w)=\varnothing$.
\end{remark}

\begin{remark}
Randomization therefore marks a boundary of CT-000. Deterministic admissibility
and robust almost-sure admissibility are both governed by the literal
common intersection $C(o)$. Pointwise almost-sure admissibility over
uncountable observation classes is not: a probability measure can give
measure one to each member of an uncountable family whose total intersection
is empty.

Randomization can also matter under weaker objectives. If failure with
positive probability is permitted, one may instead optimize
\[
  \sup_{\rho}
  \inf_{w\in[o]}
  \rho(o)\bigl(\admissible(w)\bigr).
\]
That is a different decision problem. CT-000 concerns guaranteed
admissibility and therefore corresponds to the value $1$ rather than to
maximization of expected or worst-case success probability.
\end{remark}

\subsection{Policy equivalence}

The existence theorem determines whether some globally admissible policy
exists, but such a policy need not be unique.

For each realized observation $o$, define
\[
  C(o)
  =
  \bigcap_{w\in[o]}\admissible(w).
\]

\begin{proposition}[Policy space]
For policies restricted to realized observations, the set of globally
admissible deterministic policies is naturally identified with
\[
  \prod_{o\in O(\worlds)} C(o).
\]
\end{proposition}

\begin{proof}
A policy is globally admissible precisely when, independently for every
realized observation $o$,
\[
  \pi(o)\in C(o).
\]
Thus specifying a globally admissible policy is equivalent to selecting one
element of each factor $C(o)$, which is exactly an element of the displayed
Cartesian product.
\end{proof}

\begin{corollary}[Uniqueness]
Suppose $C(o)\neq\varnothing$ for every $o\in O(\worlds)$, so that a globally
admissible realized-observation policy exists. Then it is unique if and only
if
\[
  |C(o)|=1
\]
for every $o\in O(\worlds)$.
\end{corollary}

\begin{remark}
The nonemptiness hypothesis is necessary. If some $C(o)$ is empty and another
has two elements, there are no globally admissible policies, so ``all
admissible policies coincide'' holds vacuously while $|C(o)|=1$ fails for
some $o$.
\end{remark}

\begin{corollary}[Finite policy count]
If $O(\worlds)$ and all $C(o)$ are finite, then the number of globally
admissible realized-observation policies is
\[
  \prod_{o\in O(\worlds)} |C(o)|.
\]
\end{corollary}

\begin{remark}
CT-000 is therefore equivalently a nonemptiness theorem for a product of
common admissible-action sets. The choice issue in the forward direction is
the set-theoretic issue of selecting an element of this product.
\end{remark}


\subsection{Manifest correspondence}

Formalizations, implementations, or tests may contain a \texttt{MANIFEST:}
header listing the claim identifiers to which they correspond.

Manifest membership asserts only correspondence. It does not assert that an
artifact proves every interpretation associated with the identifier.

The registry imposes the following consistency condition:

\begin{quote}
Every \texttt{CT-\#\#\#} identifier declared by a manifest must exist in the
canonical claim registry.
\end{quote}

The converse is intentionally false: a mathematical or expository claim need
not have an implementation artifact.

Revision identity, timestamps, authorship information, and content hashes are
provided by Git and the repository hosting system. The claim registry does
not define a second provenance mechanism.

\subsection{Scope of formal verification}

A proof-assistant formalization should state exactly which mathematical
proposition it verifies.

For CT-000, a direct formalization requires only sets, functions,
observation classes, admissible-action sets, and a suitable selection
principle. It need not formalize CT-001--CT-016.

In particular, formal verification of CT-000 does not establish:

\begin{itemize}
  \item the qualitative dynamical claim CT-004;
  \item a universal dimensionality result for CT-007;
  \item a literal sheaf or Transformer interpretation for CT-008;
  \item the field-model assumptions of CT-010--CT-014; or
  \item a particular software architecture for CT-016.
\end{itemize}

Those claims require their own definitions and assumptions if they are to be
formalized.

Conversely, an implementation satisfying CT-016 does not constitute a proof
of CT-000. The connection is structural: both concern preservation of
distinctions relevant to admissible action.

\subsection{Compatibility}

Adoption of this claim registry does not require existing formalizations,
experiments, or implementations to be rewritten.

An artifact may correspond to a subset of the claims. It should identify that
subset explicitly rather than being described as an implementation of the
entire theory.

Likewise, adoption of the registry does not authorize changes to established
serialization formats, transition semantics, replay behavior, or experimental
fixtures. Those are independent compatibility constraints.

The purpose of the identifiers is therefore limited and technical: they
provide stable references between mathematical statements and artifacts
without collapsing those artifacts into a single evidential category.

% =====================================================================
\section{Epistemic Overlap}
%% CT-000

\begin{definition}[Admissible-action structure]
Let $\worlds$ be a set of worlds, $\obs$ a set of observations, and
$\acts$ a nonempty set of actions. Let
\[
  O : \worlds \to \obs
\]
be an observation map and let
\[
  \admissible : \worlds \to \mathcal{P}(\acts)
\]
assign to each world its set of admissible actions.

For $o \in O(\worlds)$, define
\[
  [o]
  =
  \{w \in \worlds : O(w)=o\}.
\]
\end{definition}

\begin{definition}[Observation-based policy]
A deterministic observation-based policy is a function
\[
  \pi : \obs \to \acts.
\]
It is \emph{globally admissible} when
\[
  \pi(O(w)) \in \admissible(w)
  \qquad
  \text{for every }w\in\worlds.
\]
\end{definition}

\begin{theorem}[Epistemic Overlap --- CT-000]
Assume a selection can be made from each nonempty common admissible-action
set. Then a globally admissible observation-based policy exists if and only if
\[
  \bigcap_{w\in[o]}\admissible(w)
  \neq\varnothing
  \qquad
  \text{for every }o\in O(\worlds).
\]
\end{theorem}

\begin{proof}
Suppose first that every observation class has a nonempty common
admissible-action set. For each $o\in O(\worlds)$ choose
\[
  a_o
  \in
  \bigcap_{w\in[o]}\admissible(w).
\]
Because $\acts$ is nonempty, extend this assignment arbitrarily to
observations in $\obs\setminus O(\worlds)$ and thereby define
\[
  \pi:\obs\to\acts.
\]
For every $w\in\worlds$,
\[
  \pi(O(w))
  \in
  \admissible(w),
\]
so $\pi$ is globally admissible.

Conversely, suppose a globally admissible policy $\pi$ exists. Fix
$o\in O(\worlds)$. For every $w\in[o]$,
\[
  \pi(o)
  =
  \pi(O(w))
  \in
  \admissible(w).
\]
Hence
\[
  \pi(o)
  \in
  \bigcap_{w\in[o]}\admissible(w),
\]
so the intersection is nonempty.
\end{proof}

\begin{corollary}[No Safe Action]
If some realized observation satisfies
\[
  \bigcap_{w\in[o]}\admissible(w)
  =
  \varnothing,
\]
then no deterministic observation-based policy can guarantee an admissible
action in every world.
\end{corollary}

\begin{corollary}[Two-world obstruction]
If there exist $w,w'\in\worlds$ such that
\[
  O(w)=O(w')
\]
and
\[
  \admissible(w)\cap\admissible(w')
  =
  \varnothing,
\]
then no globally admissible observation-based policy exists.
\end{corollary}

\begin{remark}
The two-world obstruction is sufficient but not necessary. With three or more
worlds, every pair of admissible-action sets may intersect while their total
intersection is empty. The full observation-class intersection in CT-000 is
therefore the appropriate condition.
\end{remark}

\begin{remark}
CT-000 is deliberately elementary. Its role is to isolate an information
constraint: per-world feasibility does not imply information-relative
controllability.

The subsequent claims investigate particular ways in which admissibility,
information, local consistency, and reconstruction can acquire additional
structure. They are applications and extensions of this distinction, not
premises required for CT-000.
\end{remark}

% =====================================================================
\section{Admissible Histories}
%% CT-001 -- CT-005

Let a history be a finite sequence
\[
  h=(x_0,x_1,\ldots,x_T)
\]
and let $\histories$ denote the set of histories permitted by the transition
structure of a system. Let
\[
  \valuation:\histories\to\mathbb{R}
\]
be a valuation defined on those histories.

\begin{proposition}[Admissibility and valuation --- CT-001]
If admissibility is defined solely by membership in $\histories$, then
membership is independent of the valuation:
\[
  h\in\histories
\]
does not depend on the value of $\valuation(h)$.
\end{proposition}

\begin{proof}
This follows directly from the definition. The transition structure determines
membership in $\histories$, whereas $\valuation$ is a function defined on the
resulting histories. Changing $\valuation$ without changing the transition
structure therefore leaves $\histories$ unchanged.
\end{proof}

\begin{remark}
The proposition distinguishes structural possibility from preference. An
admissible history may be undesirable, unrewarded, or unintended.
\end{remark}

\noindent\emph{Source:} \texttt{Admissible Histories.tex}, \S3.

\begin{proposition}[Prefix closure --- CT-002]
Assume $\histories$ is prefix-closed. Then
\[
  (x_0,\ldots,x_T)\in\histories
  \quad\Longrightarrow\quad
  (x_0,\ldots,x_t)\in\histories
\]
for every $0\leq t\leq T$.
\end{proposition}

\begin{remark}
Prefix closure is an assumption on the admissible-history model, not a
property of arbitrary sets of trajectories.
\end{remark}

\noindent\emph{Source:} \texttt{Admissible Histories.tex}, \S2.

\begin{proposition}[Unique continuation --- CT-003]
Let $p=h_{\leq t}$ be an admissible prefix. Suppose $h$ is the unique
admissible history extending $p$ through a specified horizon. If another
extension $h'$ shares the prefix $p$ but differs from $h$ after $t$, then
\[
  h'\notin\histories.
\]
\end{proposition}

\begin{proof}
By hypothesis, $h$ is the unique member of $\histories$ extending $p$ through
the relevant horizon. Any distinct extension therefore lies outside
$\histories$.
\end{proof}

\begin{remark}
A prefix of this kind can be interpreted as a commitment because it reduces
the admissible continuation set to a singleton. More generally, commitment
can be represented by a reduction in the continuation set without requiring
uniqueness.
\end{remark}

\noindent\emph{Source:} \texttt{Admissible Histories.tex}, \S2.

\begin{claim}[Persistence of disfavored admissible trajectories --- CT-004]
In models with sparse connectivity, competitive dynamics, and incremental
weight adjustment, unfavorable valuation need not eliminate all structurally
admissible trajectories leading to unfavorable outcomes.
\end{claim}

\begin{remark}
The stronger assertion that every sufficiently complex system must retain such
trajectories requires a specified dynamical and weight-space model. CT-004 is
therefore treated here as a qualitative hypothesis rather than a general
theorem.
\end{remark}

\noindent\emph{Source:} \texttt{Admissible Histories.tex}, \S4.

\begin{proposition}[Pruning intervention --- CT-005]
Let
\[
  D:\{\text{prefixes}\}\to\{0,1\}
\]
be a detection rule and define
\[
  \histories'
  =
  \{h\in\histories : D(h_{\leq t})=0\}.
\]
Then pruning changes the admissible-history set by excluding continuations. It
does not retroactively modify the prefix on which the decision was made.
\end{proposition}

Suppose additionally that $\histories$ is finite and that $\histories'$ is
nonempty. Define
\[
  \operatorname{Acc}(\mathcal{S})
  =
  \frac{
    |\{h\in\mathcal{S}:\valuation(h)>0\}|
  }{
    |\mathcal{S}|
  },
  \qquad
  P=|\{h\in\histories:\valuation(h)>0\}|,
  \quad
  N=|\{h\in\histories:\valuation(h)\leq0\}|,
\]
and assume $P,N>0$. Let $r_P$ and $r_N$ be the numbers of positive- and
nonpositive-valuation histories removed by $D$. Then
\[
  \operatorname{Acc}(\histories)=\frac{P}{P+N},
  \qquad
  \operatorname{Acc}(\histories')=\frac{P-r_P}{P+N-r_P-r_N},
\]
and
\[
  \operatorname{Acc}(\histories')>\operatorname{Acc}(\histories)
  \iff
  \frac{r_N}{N}>\frac{r_P}{P}.
\]
Indeed, since both denominators are positive, cross-multiplication gives
$(P-r_P)(P+N)>P(P+N-r_P-r_N)$, which simplifies to $P\,r_N>N\,r_P$.
Pruning therefore increases the positive-history proportion precisely when the
proportional removal rate among nonpositive histories exceeds that among
positive histories.

The structural content of CT-005 is independent of this inequality: pruning
operates by restricting continuation, not by modifying the prefix on which the
decision was made.

\noindent\emph{Source:} \texttt{Admissible Histories.tex}, \S5.

% =====================================================================
\section{Local and Global Geometry}
%% CT-006 -- CT-009

Let $M\subset\mathbb{R}^n$ be a smooth embedded $d$-dimensional submanifold.

\begin{theorem}[Tangent-normal decomposition --- CT-006]
For every $x\in M$,
\[
  \mathbb{R}^n
  =
  T_xM\oplus N_xM,
\]
where
\[
  N_xM=(T_xM)^\perp.
\]
Consequently every ambient vector $v\in\mathbb{R}^n$ admits a unique
decomposition
\[
  v=v_T+v_N
\]
with $v_T\in T_xM$ and $v_N\in N_xM$.
\end{theorem}

\begin{remark}
The decomposition is a standard result of differential geometry. The
additional interpretation used here identifies tangent motion with locally
structure-preserving transformation and normal motion with departure from a
modeled constraint manifold. That interpretation is a modeling choice rather
than part of the geometric theorem.
\end{remark}

\noindent\emph{Source:} \texttt{Never-Predict-Noise.tex}, tangent-normal
decomposition; see also
\texttt{physics/cosmology/admissible-geometry.tex}.

\begin{claim}[Off-manifold generation --- CT-007]
Within the geometric model above, persistent nonzero normal components can
move generated trajectories away from the modeled admissible manifold.

Under additional regularity, rank, and measure assumptions, such departures
may enlarge the effective support of generated outputs relative to the
original $d$-dimensional manifold.
\end{claim}

\begin{remark}
A nonzero normal component alone does not imply, without further hypotheses,
that the output support must have Hausdorff dimension greater than $d$.
Accordingly, no universal dimensionality-increase theorem is asserted here.

Likewise, ``off-manifold'' is used as a mathematical model of a class of
generation errors. It is not proposed as a universal definition of
hallucination.
\end{remark}

\noindent\emph{Source:} \texttt{Never-Predict-Noise.tex}, off-manifold
analysis.

\begin{claim}[Local-chart coherence --- CT-008]
Suppose $M$ is covered by local coordinate neighborhoods $U_i$. A global
description assembled from local representations requires compatibility on
their overlaps.

This supplies a model for computational systems in which locally produced
representations must remain mutually consistent where their effective domains
overlap.
\end{claim}

A related constrained-optimization model uses projected updates of the form
\[
  x_{t+1}
  =
  x_t
  -
  \eta\,
  \Pi_{T_{x_t}M}
  \nabla\mathcal{L}(x_t),
\]
where $\Pi_{T_{x_t}M}$ is the orthogonal projection onto the tangent space.

\begin{remark}
These constructions provide mathematical analogies for local representation
and constrained update. They do not imply that an arbitrary attention
mechanism literally implements a sheaf or that an arbitrary Transformer layer
implements the displayed tangent-flow equation.
\end{remark}

\noindent\emph{Source:} \texttt{Never-Predict-Noise.tex}, local-chart and
tangent-flow discussion.

\begin{claim}[Local versus global coherence --- CT-009]
Local satisfaction of an admissibility criterion does not in general imply
global satisfaction of the corresponding composed criterion.

In geometric approximations, individually acceptable local steps may
accumulate projection, discretization, curvature, or chart-transition error
over a sufficiently long trajectory. A trajectory can therefore depart
globally from the intended region even though each individual step satisfies
the chosen local approximation.
\end{claim}

\begin{remark}
The connection with CT-000 is structural. CT-000 requires compatibility across
all worlds represented by an observation. CT-009 emphasizes the analogous
fact that locally acceptable choices require an additional compatibility
condition before they can be composed into a globally acceptable trajectory.
\end{remark}

% =====================================================================
\section{Visibility and Bounded Correction}
%% CT-010 -- CT-014

The following results belong to a particular field model. They are conditional
on its assumptions and should not be interpreted as universal laws of
communication or visibility.

Let $\Phi(x,t)$ denote visibility, $\mathbf{v}(x,t)$ an interpretive-flow
field, $S(x,t)$ a disorder variable, and $\mathcal{C}(x,t)$ a corrective term
over a finite-measure domain $\manifold$. Let
\[
  \alpha>0,\qquad \beta>0,\qquad c_1>0,
\]
and let protective capacity satisfy
\[
  0<P(x,t)\leq P_{\max}<\infty.
\]

We consider the following assumptions.

\begin{description}
  \item[A1. Entropy production.]
  \[
    \frac{\partial S}{\partial t}
    \geq
    \alpha|\mathbf{v}|^2
    -
    \beta\mathcal{C},
    \qquad
    \alpha>0.
  \]

  \item[A2. Flow coupling.]
  \[
    |\mathbf{v}|
    \geq
    c_1|\nabla\Phi|.
  \]

  \item[A3. Increasing interaction density along a scaling family.]
  Let $\{\Phi_\lambda\}_{\lambda\in\Lambda}$ be the scaling family under
  consideration, with associated fields $\mathbf{v}_\lambda$ and $S_\lambda$
  satisfying A1, A2 and A4 with constants $\alpha,\beta,c_1,C_{\max}$
  independent of $\lambda$. Along a directed regime $\lambda\to\lambda_*$,
  \[
    \int_{\manifold}|\nabla\Phi_\lambda|^2\,dx
    \longrightarrow\infty.
  \]
  This scaling limit is distinct from temporal evolution in $t$.

  \item[A4. Bounded correction.]
  \[
    \mathcal{C}(x,t)
    \leq
    C_{\max}<\infty.
  \]
\end{description}

\begin{proposition}[Bounded correction under increasing visibility --- CT-010]
Under A1--A4, the positive production term can become arbitrarily large
relative to the bounded correction term along the assumed scaling regime.
Consequently, a fixed bounded protective envelope cannot in general be
guaranteed for arbitrarily increasing visibility.
\end{proposition}

\begin{proof}
Fix a member $\Phi_\lambda$ of the scaling family and write $\Phi$,
$\mathbf{v}$, $S$ for $\Phi_\lambda$, $\mathbf{v}_\lambda$, $S_\lambda$. By A2,
\[
  |\mathbf{v}|^2
  \geq
  c_1^2|\nabla\Phi|^2.
\]
Substitution into A1 gives
\[
  \frac{\partial S}{\partial t}
  \geq
  \alpha c_1^2|\nabla\Phi|^2
  -
  \beta\mathcal{C}.
\]
Integrating over $\manifold$ yields
\[
  \frac{d}{dt}
  \int_{\manifold}S\,dx
  \geq
  \alpha c_1^2
  \int_{\manifold}|\nabla\Phi|^2\,dx
  -
  \beta
  \int_{\manifold}\mathcal{C}\,dx.
\]
When $\manifold$ has finite measure, A4 bounds the second integral by
\[
  C_{\max}|\manifold|.
\]
The constants $\alpha,\beta,c_1,C_{\max}$ and $|\manifold|$ do not depend on
$\lambda$, so by A3 the lower bound
\[
  \alpha c_1^2\int_{\manifold}|\nabla\Phi_\lambda|^2\,dx
  -\beta C_{\max}|\manifold|
\]
tends to infinity along the scaling regime. Thus bounded correction cannot
uniformly offset the modeled production term.

Additional assumptions relating integrated growth of $S$ to the pointwise
protective bound $P$ determine the precise threshold at which
$S(x,t)>P(x,t)$ must occur.
\end{proof}

\begin{corollary}[Finite safe regime --- CT-011]
If the additional assumptions of the model yield a monotone relation between
visibility and accumulated disorder, while protective capacity remains
bounded, then the regime in which
\[
  S(x,t)\leq P(x,t)
\]
can extend only to a finite range of the corresponding visibility parameter.
\end{corollary}

\begin{remark}
The result does not define a universal numerical visibility threshold. Any
threshold depends on the parameters, domain, initial conditions, and scaling
relations of the model.
\end{remark}

\begin{claim}[Risk growth --- CT-012]
If the visibility model is supplemented by scaling relations under which the
disorder-producing term grows faster than the available corrective term, then
marginal modeled risk may increase with visibility over that regime.
\end{claim}

\begin{remark}
Superlinear risk growth is therefore model-dependent. It is not asserted as a
general property of visibility, communication, popularity, or information
flow.
\end{remark}

\begin{claim}[Local stress criterion --- CT-013]
A field formulation may decompose local semantic stress into a
coherence-preserving contribution $\Theta_\phi$ and a
contradiction-producing contribution $\Theta_{S,\Phi}$.

Within such a model, a sufficient local stability condition can take the form
\[
  \left|
    \nabla\cdot\Theta_\phi
  \right|
  \geq
  \left|
    \nabla\cdot\Theta_{S,\Phi}
  \right|.
\]
\end{claim}

\begin{remark}
The displayed inequality is a criterion internal to the field model. It is
not asserted to be a general physical stress law.
\end{remark}

\begin{claim}[Scaling formulation --- CT-014]
Suppose a model has scale-dependent visibility
\[
  \Phi_\ell
  \sim
  \ell^{\Delta_\Phi},
  \qquad
  \Delta_\Phi>0,
\]
and suppose the associated disorder term grows more rapidly with $\ell$ than
the stabilizing contribution. Then a dimensionless stability ratio may
decrease as scale increases.
\end{claim}

\begin{remark}
Renormalization terminology is used here as a scaling framework. The claim
does not require interpreting visibility as a literal physical
renormalization-group operator. Its content depends on the specified scaling
relations.
\end{remark}

\noindent\emph{Source for CT-010--CT-014:}
\texttt{Against\_Indiscriminate\_Visibility.tex}.

% =====================================================================
\section{Decision-Relevant History and Durable Adjudication}
%% CT-015 -- CT-016

\begin{definition}[Decision-relevant factorization]
Let $H$ be a space of complete histories and suppose
\[
  q:H\to M
\]
maps histories to a reduced decision state. The admissibility predicate
\emph{factors through} $q$ when there exists
\[
  \admissible_M:M\to\mathcal{P}(\acts)
\]
such that
\[
  \admissible(h)
  =
  \admissible_M(q(h))
\]
for every relevant history $h$.
\end{definition}

\begin{proposition}[Decision-relevant history --- CT-015]
When admissibility factors through $q$, complete history is unnecessary for
the admissibility decision itself: histories $h$ and $h'$ satisfying
\[
  q(h)=q(h')
\]
have the same admissible-action set.
\end{proposition}

\begin{proof}
By factorization,
\[
  \admissible(h)
  =
  \admissible_M(q(h))
  =
  \admissible_M(q(h'))
  =
  \admissible(h').
\]
\end{proof}

\begin{remark}
The proposition establishes sufficiency, not minimality. Showing that $M$ is
a \emph{minimal} decision-relevant representation requires a separate
minimality criterion.

In an operational system, $M$ might contain the last attested state, a
monotone sequence position, and any other information required to determine
legal continuation.
\end{remark}

\noindent\emph{Source:} \texttt{Admissible Histories.tex}, \S\S2--3.

\begin{claim}[Durable adjudication --- CT-016]
Suppose an adjudication changes the set of legal future continuations. If
authoritative state must be reconstructible after restart, information
necessary to reproduce that adjudication must survive reconstruction.

In particular, when a refusal affects subsequent admissibility, recording only
successful state transitions is insufficient. The refusal itself must be
represented in whatever durable state is used to reconstruct future legal
continuations.
\end{claim}

\begin{remark}
One implementation strategy is an append-only adjudication log in which both
admissions and refusals receive monotone positions. Submitted client sequence
information can remain distinct from authoritative committed sequence
information, allowing drift or replay to be detected.

This is an implementation principle derived from the continuation model, not
an independent mathematical theorem.
\end{remark}

\begin{proposition}[Reconstruction invariant]
Let $L$ be a durable adjudication log and let
\[
  F(s_0,L)
\]
reconstruct authoritative state from an initial state $s_0$. If future
admissibility depends on an adjudication $e$, then any durable representation
claimed to preserve authoritative behavior across restart must retain enough
information about $e$ that reconstruction yields the same future
admissibility relation.
\end{proposition}

\begin{proof}
Otherwise there exist two histories differing in the relevant adjudication
but reconstructing to states indistinguishable to the future admissibility
procedure. If the adjudication changes legal continuation, those reconstructed
states require different admissible-action sets. The reconstruction has
therefore erased an admissibility-relevant distinction, producing exactly the
kind of information loss identified by CT-000.
\end{proof}

% =====================================================================
\section{Logical Status of the Claims}

The claim identifiers provide stable references, but they do not imply that
all claims have the same epistemic or mathematical status.

CT-000 is the central abstract theorem. CT-001--CT-005 describe the
admissible-history model; several follow from explicit structural assumptions.
CT-004 remains qualitative.

CT-006 contains a standard geometric theorem together with an explicitly
separate modeling interpretation. CT-007--CT-009 concern the use of geometric
structures as models of local and global computational consistency and require
the hypotheses stated in their respective formulations.

CT-010--CT-014 belong to a conditional visibility-and-correction model. Their
applicability depends on the definitions, field equations, and scaling
assumptions of that model.

CT-015 is a factorization result. CT-016 is principally an implementation
constraint motivated by the continuation model.

A machine-checked formalization of one of these claims establishes only the
formal statement encoded there. An implementation corresponding to a claim
demonstrates implementation behavior, not a proof of the unrestricted
mathematical proposition. Conversely, a mathematical theorem does not by
itself establish that a particular implementation satisfies its hypotheses.

This separation is intentional. The claim registry provides correspondence
without collapsing proof, formalization, modeling, and implementation into a
single category.

% =====================================================================
\section*{Claim Index}

\begin{center}
\small
\begin{tabular}{@{}p{1.5cm}p{6.4cm}p{5.2cm}@{}}
\toprule
Identifier & Description & Status \\
\midrule
CT-000 & Epistemic Overlap & proved \\
CT-001 & Admissibility and valuation & established within the stated model \\
CT-002 & Prefix closure & established under the prefix-closure assumption \\
CT-003 & Unique continuation & established within the stated model \\
CT-004 & Persistence of disfavored admissible trajectories & qualitative \\
CT-005 & Pruning intervention & established within the stated model \\
CT-006 & Tangent-normal decomposition & standard result plus modeling correspondence \\
CT-007 & Off-manifold generation & conditional \\
CT-008 & Local-chart coherence & modeling correspondence \\
CT-009 & Local versus global coherence & general structural observation; geometric form model-dependent \\
CT-010 & Bounded correction under increasing visibility & conditional \\
CT-011 & Finite safe regime & conditional \\
CT-012 & Risk growth & model-dependent \\
CT-013 & Local stress criterion & model-dependent \\
CT-014 & Scaling formulation & model-dependent \\
CT-015 & Decision-relevant history & proved as a factorization result \\
CT-016 & Durable adjudication & implementation principle \\
\bottomrule
\end{tabular}
\end{center}

\bigskip

\noindent
The canonical definitions and status of these identifiers are maintained in
\texttt{SPEC.md}.

\end{document}

